\section{Introduction}
\label{sec:introduction}

The deployment of Large Language Models (LLMs) across multilingual societies has outpaced our ability to measure and guarantee their factual trustworthiness. While state-of-the-art foundation models demonstrate remarkable fluency across diverse natural languages, their internal knowledge representations exhibit severe structural asymmetries~\citep{ahuja2023mega,bang2023multitask}. In particular, when generating text in non-Latin, lower-resource, or regional languages, LLMs frequently produce outputs that are grammatically coherent yet factually fabricated---a phenomenon broadly formalized as \emph{hallucination}~\citep{maynez2020faithfulness,ji2023survey,tonmoy2024comprehensive}.

Despite substantial recent progress in factuality benchmarking---such as TruthfulQA~\citep{lin2022truthfulqa}, FActScore~\citep{min2023factscore}, and HaluEval~\citep{li2023halueval}---the vast majority of empirical investigations remain disproportionately confined to English. Evaluating factuality in low-resource and non-Western languages presents acute conceptual and methodological hurdles that standard English benchmarks fail to capture.

First, standard subword tokenizers (e.g., Byte-Pair Encoding and SentencePiece), trained predominantly on Latin-script corpora, heavily fragment non-Latin orthographies into sub-character byte sequences~\citep{hasan2023banglatokenizer}. In Bengali, complex consonant conjuncts (\foreignlanguage{bengali}{যুক্তাক্ষর}) and vowel diacritics (\foreignlanguage{bengali}{কার}) suffer from extreme fragmentation ratios ($0.51$--$0.88$ compression relative to English), requiring up to double the token budget and diluting semantic attention across autoregressive decoding steps.

Second, digital communication in South Asia is characterized by pervasive diglossia and informal Romanized transliteration. Hundreds of millions of Bengali speakers routinely converse in \emph{Banglish}---phonetically transcribed Bengali rendered in the Latin alphabet (e.g., \emph{"amader shongbidhan kobe pass hoyesilo"} for \foreignlanguage{bengali}{"আমাদের সংবিধান কবে পাস হয়েছিল?"}). Because Banglish lacks standardized orthography, identical morphemes undergo extreme phonetic variance (e.g., \emph{kisu} versus \emph{kichu}, \emph{valo} versus \emph{bhalo}), introducing severe out-of-vocabulary dispersion into multilingual embedding spaces~\citep{saha2024robustness,chowdhury20245dialects}.

Third, localized societal applications in public health, clinical medicine, and statutory jurisprudence carry severe real-world hazards when subjected to ungrounded generation~\citep{pal2023medhalt,guha2023legalbench}. An LLM providing hallucinated pharmacotherapy advice or inventing non-existent constitutional provisions directly endangers human health or misinforms civic legal rights.

Finally, while Retrieval-Augmented Generation (RAG)~\citep{lewis2020rag} is commonly deployed to supply external context, multilingual dense retrievers frequently surface noisy or off-target distractor passages~\citep{cuconasu2024power,park2024robust}. When presented with ungrounded context, models must possess the epistemic capability to \emph{selectively abstain}~\citep{cole2023selectively,yin2023large} rather than confabulating plausible-sounding misinformation.

To systematically investigate these phenomena, we introduce \textbf{Matibhrom (\foreignlanguage{bengali}{মতিভ্রম})}---derived from the Bengali word for cognitive illusion or aberrant perception---the first fine-grained, multi-condition hallucination benchmark constructed specifically for the Bengali language. Matibhrom establishes an end-to-end experimental framework designed around five foundational research questions: investigating the factual impact of informal Latin transliteration compared to native script (RQ1), disentangling parametric memory failure from extractive reading comprehension deficits (RQ2), measuring susceptibility to realistic dense retrieval distractors in RAG pipelines (RQ3), evaluating whether models can reliably detect evidence shortfalls and abstain under explicit negative constraints (RQ4), and analyzing how parameter scaling governs the emergence of calibrated abstention in low-resource settings (RQ5).

In addressing these questions, this work delivers four principal contributions. First, we construct and release the Matibhrom benchmark, comprising 600 expert-verified base items (yielding 1,200 matched prompts in native Standard Bangla and Banglish) grounded across a frozen corpus of 1,800 passages in General, Medical, and Legal domains, partitioned into document-isolated, leakage-free splits. Second, we design a standardized 4-condition grounding evaluation protocol spanning Closed-book ($C_0$), Oracle Grounding ($C_1$), Dense RAG with Distractors ($C_2$), and Selective Abstention ($C_3$), indexed via dense multilingual representations (\texttt{BAAI/bge-m3}). Third, across 4,560 evaluated test generations evaluating Qwen 2.5 1.7B and Gemma 3 4B, we report 95\% Wilson score confidence intervals and paired McNemar significance tests ($\chi^2$) with continuity corrections, providing rigorous statistical proof for every claim. Most notably, we discover that parameter scaling from 1.7B to 4B parameters unlocks emergent selective abstention under noisy retrieval ($\chi^2 = 51.191, p = 8.38 \times 10^{-13}$), slashing hallucinations from 37.5\% to 3.3\% with 92.8\% abstention. Finally, we demonstrate that automated zero-shot LLM judges exhibit unacceptable label variance ($\kappa \le 0.34$) on low-resource Bengali, establishing the methodological imperative for deterministic gold-adjudication frameworks in multilingual evaluation.
